% simulation.tex -- Methodology 3.6: the experimental protocol (rewritten 12 Oct 2026 to the thesis plan of 11 Oct,
% "3.6 Method"). The rig twin's description moved to system_overview.tex (3.1) with its label sec:simulation.
% Sources: tools/run_experiment.py, results/registry.csv, docs/experiments/*.md (the cards), the critic and reviewer
% agents of the working rules, docs/freeze_2026-10-14.md (the signed rules). Counts (runs, tests) are left in
% brackets to be filled at the results freeze.

\subsubsection{Protocol}
\label{sec:protocol}

\paragraph{Experiment cards, frozen bars and one code tree.} A result that supports a claim, a thesis table or a
go/no-go for the rig, is flown under an experiment card written before the runs: the question, the one variable, the
baseline, the hypothesis with its sign, the pass and fail numbers and the repeats. The bars are frozen before the runs,
and a bar found wrong applies only to later runs. Every arm of a claim table is flown on one code tree. Exploratory runs,
which find a failure rather than prove a fix, carry no card and are marked as such; every run, voided ones included, is
recorded in a register. The rig flies on the tree of its lab day, so each rig row carries its own tree and its twin pair
is flown on that tree before the flight.

\paragraph{Checks before and after the runs.} Two automated checks sit around every claim, and neither is an
independent review. Before a card runs, a critic reads it against the code and the register and challenges the
variable, the baseline and the bars; its challenge and the answer are kept on the card. After the runs, a reviewer
reads the claim against the run data and the frozen bars and may pass it, fail it or send it back, and no number
reaches a thesis table without that pass. Both are large-language-model agents run against pre-frozen bars and
the logged data, and the author decides what they contest. [The disclosure of AI assistance, in the form the
University requires, stands before this section.]

\paragraph{Runner and ladder.} A twin run is a configuration file (the world, the launch arguments, a timed list of
events and the pass and fail criteria) that a runner flies headless from a clean machine, reading back the
parameters every node loaded and computing the metrics from the logs. Changes climb a ladder from unit tests to a
quick check script, headless twin runs and only then the rig. Cases are flown twice in the twin, and tables give the
number of runs; rig cells are flown as often as a lab day allows and say so.
